% --- 题目 ---
\problem[P10-4 (02-2)]{$\displaystyle \lim_{n \to \infty} \frac{1}{n} \left[ \sqrt{1+\cos \frac{\pi}{n}} + \sqrt{1+\cos \frac{2\pi}{n}} + \ldots + \sqrt{1+\cos \frac{n\pi}{n}} \right] = \underline{\hspace{4em}}. $}
\ansat{243；【真题精选】 - 考点二：数列极限的计算 - 4}

% --- 题目 ---
\problem[P10-8 (99-2)]{设 $\displaystyle f(x)$ 是区间 $\displaystyle [0, +\infty)$ 上单调减少且非负的连续函数.
\[ a_n = \displaystyle \sum_{k=1}^{n} f(k) - \int_{1}^{n} f(x) \ \mathrm{d}x \ (n = 1, 2, \dots), \]
证明数列 $\displaystyle \{a_n\}$ 的极限存在.}
\ansat{243；【真题精选】 - 考点二：数列极限的计算 - 8}

% --- 题目 ---
\probtagged[P8-变式 4.1 (96-1)]{设 $\displaystyle x_1 = 10, \ x_{n+1} = \sqrt{6+x_n} \ (n = 1, 2, \dots)$. 试证数列 $\displaystyle \{x_n\}$ 极限存在, 并求此极限.}
\ansat{241；【方法探究】 - 考点二：数列极限的计算 - 变式 4.1}

% --- 题目 ---
\problem[P8-变式 3 (98-1)]{求极限 $$\lim_{n \to \infty} \left( \frac{\sin \frac{\pi}{n}}{n+1} + \frac{\sin \frac{2\pi}{n}}{n + \frac{1}{2}} + \dots + \frac{\sin \pi}{n + \frac{1}{n}} \right).$$}
\ansat{240；【方法探究】 - 考点二：数列极限的计算 - 变式 3}

% --- 题目 ---
\problem[P10-1 (12-2)]{设 $\displaystyle a_n > 0 \ $ ($\displaystyle n=1,2,\dots$), $\displaystyle S_n = a_1 + a_2 + \dots + a_n$, 则数列 $\displaystyle \{S_n\}$ 有界是数列 $\displaystyle \{a_n\}$ 收敛的( \quad )}
\begin{tasks}(2)
\task 充分必要条件.
\task 充分非必要条件.
\task 必要非充分条件.
\task 既非充分也非必要条件.
\end{tasks}
\ansat{243；【真题精选】 - 考点二：数列极限的计算 - 1}

% --- 题目 ---
\problem[P4-5 (19-1,3)]{设 $\displaystyle a_n = \int_{0}^{1} x^n \sqrt{1 - x^2} \mathrm{d}x \ (n = 0, 1, 2, \dots)$.}
\begin{enumerate}
\item[(1)] 证明: 数列 $\displaystyle \{a_n\}$ 单调减少, 且 $\displaystyle a_n = \frac{n-1}{n+2} a_{n-2} \ (n=2, 3, \dots)$;
\item[(2)] 求 $\displaystyle \lim_{n \to \infty} \frac{a_n}{a_{n-1}}$.
\end{enumerate}
\ansat{239；【十年真题】 - 考点二：数列极限的计算 - 5}

% --- 题目 ---
\problem[P4-6 (18-1,2,3)]{设数列 $\displaystyle \{x_n\}$ 满足：$\displaystyle x_1 > 0, \ x_n \mathrm{e}^{x_{n+1}} = \mathrm{e}^{x_n} - 1 \ (n = 1, 2, \dots)$. 证明 $\displaystyle \{x_n\}$ 收敛, 并求 $\displaystyle \lim_{n \to \infty} x_n$.}
\ansat{240；【十年真题】 - 考点二：数列极限的计算 - 6}